\chapter{Apertura del Tomo II}

\section{Función del tomo}
\begin{observation}
El Tomo II recibe como punto de partida el fundamento ya fijado editorialmente en Tomo I y organiza el paso desde el núcleo estructural mínimo hacia formulaciones disciplinadas de régimen, transición y jerarquía.
\end{observation}

\begin{observation}
Su función no es refundar primitivas ni duplicar el papel del Tomo I. Su tarea consiste en tomar un dominio ya estabilizado en sentido mínimo y prepararlo para tratamiento operativo más fuerte.
\end{observation}

\section{Relación con Tomo I}
\begin{observation}
Tomo I fija, como volumen fundacional principal, el fundamento del sistema, su línea de lectura sobria y el núcleo estructural mínimo necesario para sostener el resto del corpus.
\end{observation}

\begin{observation}
Ese núcleo comprende, al menos:
\begin{enumerate}
  \item primitivas;
  \item medida mínima;
  \item involución mínima;
  \item plano carlesino;
  \item trayectoria mínima;
  \item lectura proyectada;
  \item coherencia mínima;
  \item acumulación mínima;
  \item criterio mínimo de régimen verificable.
\end{enumerate}
\end{observation}

\begin{observation}
Los materiales técnicos anteriormente separados como Base Axiomática pueden conservarse como soporte interno de referencia, pero no alteran la jerarquía pública del fundamento ya integrado en Tomo I.
\end{observation}

\begin{observation}
A partir de ese fundamento ya fijado, el Tomo II introduce formulaciones más disciplinadas para:
\begin{enumerate}
  \item transición entre niveles;
  \item criterios operativos de régimen;
  \item umbrales;
  \item articulación jerárquica;
  \item preparación de formulaciones canónicas utilizables por demostraciones, protocolos y demos.
\end{enumerate}
\end{observation}

\section{Cláusula de no duplicación}
\begin{observation}
Nada de lo que entre al Tomo II debe volver a presentar como primitivo aquello que ya quedó fijado en Tomo I, salvo por referencia disciplinada y estrictamente necesaria.
\end{observation}

\begin{observation}
Cuando un bloque del Tomo II dependa de una pieza basal, esa dependencia debe leerse como herencia explícita y no como refundación del sistema.
\end{observation}

\section{Régimen de lectura}
\begin{observation}
El Tomo II opera en un nivel intermedio: más fuerte que la formulación mínima de base, pero todavía más austero que la capa de demos, laboratorio y traducción a claims seleccionados.
\end{observation}

\begin{observation}
Su éxito editorial depende de mantener visible esa posición intermedia sin deslizarse ni hacia un regreso a los primitivos ni hacia una mezcla prematura con casos experimentales.
\end{observation}

\section{Cláusula de alcance}
\begin{observation}
Este tomo desarrolla canon operativo, transición y jerarquía. No reemplaza a Tomo I como volumen fundacional ni absorbe la función metodológica y auditora propia del Tomo III.
\end{observation}

